\chapter{Simulation and Monte Carlo Methods}
\label{ch:sim}

Simulation turns probability theory into an experimental tool. Given a source of
independent uniform random variables (Theorem~\ref{thm:iidexist}), we can generate any
distribution, estimate expectations, and measure the error of the estimate. This chapter
shows the standard methods and \emph{why} they work.

\section{Generating random variables}

\subsection{Inverse transform}

\begin{theorem}[Inverse transform sampling]\label{thm:invtransform}
Let $F$ be any CDF and define the generalised inverse
$F^{-1}(u)=\inf\{x:F(x)\ge u\}$ for $u\in(0,1)$. If $U\sim U(0,1)$ then $X=F^{-1}(U)$ has CDF $F$.
\end{theorem}
\begin{proof}
By right-continuity of $F$, $F^{-1}(u)\le x\iff u\le F(x)$. (If $u\le F(x)$ then $x$ belongs to
the set whose infimum is $F^{-1}(u)$. Conversely if $F^{-1}(u)\le x$ then by
right-continuity the infimum is attained, $F(F^{-1}(u))\ge u$, and $F$ is non-decreasing.)
Hence $\Prob(X\le x)=\Prob(U\le F(x))=F(x)$.
\end{proof}

\begin{example}
\emph{Exponential}: $F(x)=1-e^{-\lambda x}$, $F^{-1}(u)=-\ln(1-u)/\lambda$; since $1-U\sim U$, use
$X=-\ln U/\lambda$. \emph{Discrete}: with $p_1,\dots,p_m$, return the smallest $k$ with
$U\le p_1+\dots+p_k$. \emph{Geometric}: $X=\lceil\ln U/\ln(1-p)\rceil$.
\end{example}

\subsection{Rejection sampling}

Suppose the target has density $f$ that is hard to invert, but there is a proposal
density $g$ we can sample and a constant $M$ with $f(x)\le Mg(x)$ for all $x$.

\begin{theorem}[Rejection sampling]\label{thm:reject}
Repeat: draw $Y\sim g$ and $U\sim U(0,1)$ independently; accept $Y$ if $U\le f(Y)/(Mg(Y))$.
The accepted value has density $f$, and the number of trials until acceptance is
geometric with mean $M$.
\end{theorem}
\begin{proof}
One trial accepts with probability
$\int g(y)\frac{f(y)}{Mg(y)}dy=\frac1M$. For a set $A$,
$\Prob(Y\in A,\text{accept})=\int_Ag(y)\frac{f(y)}{Mg(y)}dy=\frac1M\int_Af$. Dividing,
$\Prob(Y\in A\given\text{accept})=\int_Af$. Trials are independent, so the number of trials is
geometric$(1/M)$.
\end{proof}

\begin{nuance}{why the constant matters}
The expected cost is $M$ proposals per sample, and $M\ge1$ with equality iff $f=g$. In
dimension $d$, natural proposals have $M$ growing exponentially in $d$ (e.g.\ sampling a
ball by rejection from its bounding cube: acceptance probability $\to0$ super-exponentially),
the ``curse of dimensionality''.
\end{nuance}

\subsection{Special methods}

\begin{itemize}
\item \textbf{Normal}: Box--Muller (Chapter~\ref{ch:jointcont}), $Z_1=\sqrt{-2\ln U_1}\cos2\pi U_2$,
$Z_2=\sqrt{-2\ln U_1}\sin2\pi U_2$.
\item \textbf{Multivariate normal}: if $\Sigma=LL^\top$ (Cholesky) and $Z$ has i.i.d.\ $N(0,1)$ entries then
$\mu+LZ\sim N(\mu,\Sigma)$ since $\Cov(LZ)=LL^\top$.
\item \textbf{Poisson process}: cumulate i.i.d.\ exponential gaps (Chapter~\ref{ch:poisson}).
\item \textbf{Non-homogeneous Poisson process by thinning}: to simulate rate $\lambda(t)\le\bar\lambda$,
generate a homogeneous process at rate $\bar\lambda$ and keep each point at time $t$
independently with probability $\lambda(t)/\bar\lambda$. This is correct by the thinning
theorem (Chapter~\ref{ch:poisson}), extended to time-dependent retention.
\item \textbf{Random permutation} (Fisher--Yates): for $i=n,\dots,2$ swap position $i$ with a uniform
position in $\{1,\dots,i\}$. It yields every permutation with probability $1/n!$ since there are
$n\cdot(n-1)\cdots2=n!$ equally likely choice sequences, each giving a distinct result.
\item \textbf{Markov chains}: apply the transition row with inverse transform at each step.
\end{itemize}

\section{Monte Carlo estimation}

To estimate $\theta=\Exp[h(X)]$, draw $X_1,\dots,X_n$ i.i.d.\ and use
$\hat\theta_n=\frac1n\sum h(X_i)$.

\begin{theorem}[Properties of the Monte Carlo estimator]
Assume $\sigma^2=\Var(h(X))<\infty$. Then
\begin{enumerate}[label=(\roman*)]
\item $\Exp\hat\theta_n=\theta$ (unbiased) and $\Var\hat\theta_n=\sigma^2/n$;
\item $\hat\theta_n\to\theta$ almost surely (strong law);
\item $\sqrt n(\hat\theta_n-\theta)\xrightarrow{d}N(0,\sigma^2)$ (CLT), so an approximate
$95\%$ confidence interval is $\hat\theta_n\pm1.96\,s_n/\sqrt n$, where $s_n$ is the sample standard
deviation.
\end{enumerate}
\end{theorem}
\begin{proof}
Linearity and independence give (i); (ii) and (iii) are the SLLN and CLT of
Chapter~\ref{ch:limits}; $s_n\to\sigma$ a.s.\ and Slutsky's theorem justify the studentised interval.
\end{proof}

\begin{nuance}{the $1/\sqrt n$ rate is dimension free}
The error is of order $\sigma/\sqrt n$ whatever the dimension of $X$---the main advantage over grid
quadrature, whose error degrades with dimension. To gain one extra digit of accuracy requires
$100$ times more samples, so the practical lever is reducing $\sigma$.
\end{nuance}

\begin{example}[Estimating $\pi$]
With $(X,Y)$ uniform on $[-1,1]^2$, $h=\indic_{X^2+Y^2\le1}$ has mean $\pi/4$ and variance
$\tfrac\pi4(1-\tfrac\pi4)\approx0.169$. The estimator $4\hat\theta_n$ has standard error
$4\sqrt{0.169/n}\approx1.64/\sqrt n$; with $n=10^6$ expect about $3$ correct decimal places
($\approx0.0016$).
\end{example}

\section{Variance reduction}

\paragraph{Antithetic variables.} Use $U$ and $1-U$ in pairs: $\hat\theta=\frac12(h(U)+h(1-U))$. Then
$\Var\hat\theta=\tfrac14\bigl(2\sigma^2+2\Cov(h(U),h(1-U))\bigr)=\tfrac12\sigma^2(1+\rho)$ where $\rho$ is
the correlation. For monotone $h$, $h(U)$ and $h(1-U)$ are negatively correlated
(a consequence of the Chebyshev association inequality, cited), so $\rho\le0$ and the variance of
a pair is at most $\sigma^2/2$---no worse than two independent draws.

\paragraph{Control variates.} Suppose $W$ has known mean $\Exp W=w$ and is correlated with $h(X)$. For any
$c$, $h(X)-c(W-w)$ is unbiased with variance
$\sigma^2-2c\Cov(h,W)+c^2\Var W$, minimised at $c^*=\Cov(h,W)/\Var W$, with minimum
$\sigma^2(1-\rho^2)$ (the same algebra as linear regression, Chapter~\ref{ch:condexp}).

\paragraph{Importance sampling.} To estimate $\Exp_f[h(X)]=\int hf$, sample from a different density $g$
(positive wherever $hf\ne0$) and weight:
\[
\Exp_f[h(X)]=\int h(x)\frac{f(x)}{g(x)}g(x)\,dx=\Exp_g\Bigl[h(Y)\frac{f(Y)}{g(Y)}\Bigr].
\]
The variance under $g$ is $\Exp_g[h^2f^2/g^2]-\theta^2$. By Cauchy--Schwarz / Jensen, it is
minimised by $g^*\propto\abs hf$, where it equals $(\Exp_f\abs h)^2-\theta^2$ (zero for $h\ge0$,
though $g^*$ requires knowing $\theta$). In practice importance sampling is essential for rare events.

\begin{example}[Rare event]
Estimate $p=\Prob(Z>5)\approx2.87\times10^{-7}$ for $Z\sim N(0,1)$. Naive Monte Carlo needs $\sim10^{9}$
samples to see a few hits. Sample $Y=5+E$ with $E\sim\mathrm{Exp}(1)$, density $g(y)=e^{-(y-5)}$ on $y>5$.
The weight is $f(y)/g(y)=\phi(y)e^{y-5}$, and $p=\Exp_g[\phi(Y)e^{Y-5}]$, with the integrand bounded
by $\phi(5)\approx1.5\times10^{-6}$ and varying slowly over the region, giving a relative
error of order $1/\sqrt n$ with a modest constant.
\end{example}

\paragraph{Stratified sampling.} Partition the space into strata of known probability and sample
within each; the variance equals the \emph{within-stratum} variance only (by the law of total
variance, Chapter~\ref{ch:condexp}), never more than naive sampling with proportional allocation.

\section{Markov chain Monte Carlo}

Often we cannot sample from the target $\pi$ directly (e.g.\ the normalising constant is
unknown). MCMC constructs a Markov chain whose stationary distribution is $\pi$ and whose
long-run averages converge to $\Exp_\pi[h]$ (ergodic theorem, Chapter~\ref{ch:markov}).

\begin{theorem}[Metropolis--Hastings]\label{thm:mh}
Let $q(x,y)$ be a proposal kernel. From state $x$, propose $y\sim q(x,\cdot)$ and accept with
probability
\[
\alpha(x,y)=\min\Bigl(1,\frac{\pi(y)q(y,x)}{\pi(x)q(x,y)}\Bigr),
\]
otherwise stay at $x$. Then the resulting chain satisfies detailed balance with respect to $\pi$,
hence $\pi$ is stationary.
\end{theorem}
\begin{proof}
For $x\ne y$ the transition density is $P(x,y)=q(x,y)\alpha(x,y)$. Then
\[
\pi(x)q(x,y)\alpha(x,y)=\min\bigl(\pi(x)q(x,y),\ \pi(y)q(y,x)\bigr),
\]
which is symmetric in $(x,y)$, so $\pi(x)P(x,y)=\pi(y)P(y,x)$. Summing (integrating) over $x$ gives
$\pi P=\pi$ as in Chapter~\ref{ch:markov}. Only the \emph{ratio} $\pi(y)/\pi(x)$ is needed, so
$\pi$ need be known only up to a constant.
\end{proof}

\begin{nuance}{what MCMC does and does not guarantee}
Stationarity is not enough: the chain must be irreducible and aperiodic for convergence, and
mixing speed (how long until the chain forgets its start) can be exceedingly slow if the target
is multimodal. Samples are correlated, so the effective sample size is $n/(1+2\sum_k\rho_k)$,
where $\rho_k$ is the lag-$k$ autocorrelation. Initial ``burn-in'' samples are discarded in practice.
\end{nuance}

The \emph{Gibbs sampler} updates one coordinate at a time from its full conditional; it is a special case
of Metropolis--Hastings with acceptance probability $1$.

\section*{Exercises}
\addcontentsline{toc}{section}{Exercises}

\begin{enumerate}
\item Give an inverse-transform sampler for the Pareto law $F(x)=1-x^{-a}$, $x\ge1$, and for the
Cauchy law.
\item Show how to sample the standard normal by rejection from the double-exponential density
$g(x)=\tfrac12e^{-\abs x}$ and find the optimal $M$ ($M=\sqrt{2e/\pi}\approx1.32$).
\item Verify the Fisher--Yates claim by induction on $n$.
\item Prove the thinning algorithm for non-homogeneous Poisson processes by computing the number of
points in $[0,t]$.
\item Estimate $\int_0^1e^{x}dx$ with antithetic variates and compute the variance reduction factor
exactly (target $e-1$).
\item Design an importance sampler for $\Prob(S_{100}>60)$ for the sum of 100 fair-coin $\pm1$ steps
using an exponentially tilted Bernoulli.
\item Show that Gibbs sampling for the bivariate normal with correlation $\rho$ is a Markov chain with
autoregressive coefficient $\rho^2$ for each coordinate, and hence the mixing slows as $\abs\rho\to1$.
\end{enumerate}
